\subsection{非共有Segreの平方・積追加型}
\label{subsec:cubic-CY-one-segre-square-split}

\begin{proposition}\label{prop:cubic-CY-one-segre-square-split-outer}
一交換子・非共有Segreで $U/U_0=k(a^2+cd)$ とする。
固定出典\cite{AS4CYpinned}の横断外側七表示、非等方外側四族、
および等方外側の $a_0,a_1,a_2,b_0,\ldots,b_4$ と上行交差 $C_\theta$ を考える。
この範囲でH3包絡となり得るのは、非等方外側の $Z(s)$、
$s^2=-1/12$ のみである。
\end{proposition}
\begin{proof}
横断七表示では中心三次は零である。
等方外側八表示と非等方外側 $N(s),C(s,0)$ では、
中心三次が非零なら可換二元三次 $k[a,b]_3$ に含まれる。
上行交差 $C_\theta$ の中心三次は零であり、$Z(s)$ では
$s^2=-1/12$ における $k[a,b]_3+kc^3+kd^3$ だけが非零となる。

これらの主張は上流の元六関係から
\texttt{verify\_cubic\_CY\_one\_segre\_square\_split.py} で検証する。
次数2から4の商射影の分母と、中心20次小行列式の分子・分母の全因子を集める。
横断 $A_2$ と上行交差 $C_\theta$ はこれらの因子が既に許容単元である。
横断 $B_{n2}(h,\mu)$ は残余 $h=\pm1$、$\mu=-1/32$ と
$\mu=(1-h)/32$ の四直線で同じ計算を行う。
一変数表示とこれらの四直線では、
残る全既約因子の数体で元関係を再計算した。
中心が残る値では、表示した語の像の張る空間に核が含まれることを確認した。
$Z$ の指定値ではその空間が六次元でかつ核に等しい。
従って $Z$ 以外のこれらの中心三次は代数閉体上で積に分解し、整域切断を与えない。

残る二つの表示 $C(s,z)$、$z\ne0$ と $L(s,h)$ をまとめて扱う。
元六関係から、$k[a,b]$ は可換で、
$c,d$ がこの基礎環を横切る一次作用は互いに逆である。
$c$ の作用の二固有値を $q_1,q_2$ とすると、
いずれも $q_i\ne1$、$q_1\ne q_2$、$q_1q_2\ne1$ である。
確認のため $C(s,z)$ では $D_\pm=4s^2\pm4s+16z+1$ と置くと、
$c$ の作用の行列は
\[
 M=D_-^{-1}\begin{pmatrix}4s^2+16z-1&-8\\8z&4s^2+16z-1\end{pmatrix}.
\]
その行列式は $D_+/D_-$、$\det(I-M)=4/D_-$ であり、
二固有値が等しいのは $z=0$ だけである。
$D_\pm$ は上流の許容条件で非零であり、$s\ne0$ なので行列式は1でない。
$L(s,h)$ の二固有値は
\[
 \frac{2s+4h+1}{2s+4h-1},\qquad
 \frac{-2s+4h-1}{-2s+4h+1}
\]
であり、同じ条件は $s,h\ne0$ と許容分母から直ちに従う。
両表示で $dc=(q_1q_2)cd$ は基礎の二次式を法とした関係である。

$c,d$ の個数で濾過する。基礎の固有基底を $x,y$ とすれば
随伴次数環の順序語 $x^i y^j c^k d^l$ は基底である。
これは元の書換え則とAS4のHilbert値を比較して従う。
この語が中心であるには
\[
 q_1^{k-l}=q_2^{k-l}=1,\qquad q_1^{i-l}q_2^{j-l}=1
\]
が必要であり、異なる語の係数は独立に比較できる。
総次数3で $k+l>0$ の場合、$|k-l|=1$ は $q_i\ne1$ から除外される。
$|k-l|=2$ なら $q_1=q_2=-1$、$|k-l|=3$ なら相異なる
原始三乗根で積1となり、いずれも矛盾する。
$k=l=1$ の場合は $i+j=1$ で、最後の式から片方の $q_i$ が1となって矛盾する。
従って中心三次の最高濾過次数は0で、全て $k[x,y]_3$ に属する。
よってこの二表示も整域切断を与えない。
\end{proof}

\begin{proposition}\label{prop:cubic-CY-one-segre-square-split-resonance}
$s^2=-1/12$ の $Z(s)$ について、原始三乗根 $q=(2s+1)/(2s-1)$ を取る。
中心三次切断は
\[
 G=P(a,b)+C c^3+D d^3,\qquad P\in k[a,b]_3
\]
で尽くされる。H3であれば $CD\ne0$ である。
\end{proposition}
\begin{proof}
中心の記述は前命題による。包絡の関係を
\[
 [a,b]=0,\quad ca=qac,\quad cb=qbc,\quad
 da=q^{-1}ad,\quad db=q^{-1}bd,\quad dc=q^2cd+\kappa a^2,qquad \kappa\ne0
\]
と書く。
$C=D=0$ は二元三次の分解で排除される。
$C=0$ なら $c$ 軸の点対象 $P_0$ を取り、$c$ の作用を恒等に規格化する。
$G$ は点イデアルの二乗に入るため、一次化は零である。
六二次関係の線型化では、$a,b,d$ の非対角成分が次次数で全て $q$ 倍される。
従って $\dim\operatorname{Ext}^1(P_0,P_0)=3$、shift の行列式は $q^3=1$。
これは補題\ref{lem:cubic-odd-Serre-determinant}に反する。
$D=0$ も $d$ 軸で逆の三乗根を用いて同じ矛盾を得る。
\end{proof}

\begin{theorem}\label{thm:cubic-CY-one-segre-square-split-resonance-exact}
前命題の $Z(s)$ で $\nu=q\kappa/(1-q)$ と置く。
切断 $G=P(a,b)+C c^3+D d^3$ が強い内在的H3を与える必要十分条件は、
\[
 CD\ne0,\qquad \deg_z P(1,z)\ge 2,\qquad
 \Delta(z)=P(1,z)^2-4CD\nu^3\text{ が平方自由}
\]
である。特に $P=b^3$、$C=D=1$ は
H3の非skew CY包絡付き切断を与える。
\end{theorem}
\begin{proof}
$CD\ne0$ は前命題で示した。
$\deg_z P(1,z)\le1$ なら $G$ は $b$ 軸の点イデアルの二乗に入る。
この点の全tails自己一次拡大は $a,c,d$ の三成分を持ち、
$b$ を恒等としたときのshiftの固有値は $1,q^{-1},q$ である。
再び奇数次元・行列式1の障害により排除される。

以下 $m=\deg_z P(1,z)\in\{2,3\}$ とする。
$\sigma(a,b,c,d)=(a,b,qc,q^{-1}d)$ による右Zhangねじりを取る。
$\sigma^3=1$ で $G$ は固定され、$c^3,d^3$ のねじり前後の積は一致する。
ねじった包絡 $B'$ では $a,b$ は中央、残る関係は
\[
 dc=q cd+(1-q)\nu a^2
\]
である。$X=c^3,Y=d^3,w=cd-\nu a^2$ と書けば
\[
 wc=qcw,\qquad w^3=XY-\nu^3a^6 .
\]
後式は $c^3d^3$ を $cd$ の三つの $q$ 共役の積に展開して得られる。
$q$-Weylの順序語基底により $B'$ は $k[a,b,X,Y]$ 上で
$c^i d^j$ ($0\le i,j<3$) を基底とする。
従って切断 $R'=B'/(P+CX+DY)$ は $k[a,b,X]$ 上階数9の自由加群で、
$Y=-(P+CX)/D$ である。

\paragraph{整域性}
$a$ を逆にし、$z=b/a$、$x=X/a^3$ と置く。
中心分数体 $K=k(a,z,x)$ への移行では $c/a,w/a^2$ が生成元となり、
三次巡回代数の二つのパラメータは
\[
 x,\qquad -g(x,z)/D,qquad
 g(x,z)=C x^2+P(1,z)x+D\nu^3 .
\]
$\Delta(z)$ が平方自由なら、その次数は $2m$ であり、
$g(x,z)=0$ は滑らかな射影モデルの種数が $m-1$ の超楕円曲線を与える。
実際、$z$ への二重被覆は $2m$ 個の単純な有限分岐点を持ち、無限遠は二点に分かれる。
曲線上の $x$ は有限点で零にも極にもならず、
無限遠の一方では位数 $m$ の零、もう一方では位数 $m$ の極を持つ。
これは二根の積 $D\nu^3/C$ が定数で、和 $-P(1,z)/C$ が次数 $m$ であることから従う。
$m=2$ では位数から $x$ は立方にならない。
$m=3$ で $x$ が立方なら、立方根は一つの単純極だけを持ち、種数2の曲線から
$\mathbb P^1$ への次数1写像を与えて矛盾する。
この非立方性は純超越変数 $a$ を加えても保たれる。
従って $g=0$ の付値で $K(\sqrt[3]{x})/K$ は非分岐で剰余次数3である。
ノルムの付値は3の倍数だが、$-g/D$ の付値は1なのでノルムでない。
よって巡回代数は分裂せず、素次数3なので斜体である。
自由基底により $R'$ はこの斜体に入る。
次数付きねじりは斉次非零元の積と最高次数項を通じて整域性を保つから、元の $R$ も整域である。

\paragraph{二つのアフィンchartの局所次数}
$c$ を逆にした次数零部分を取る。$\xi=ac^{-1}$、$\eta=bc^{-1}$、
$v=wc^{-2}$ とすればこの代数は
\[
 k[\xi,\eta][v;\tau]/(v^3-F),\qquad
 \tau(\xi,\eta)=(q^{-1}\xi,q^{-1}\eta),\qquad
 F=-(C+P(\xi,\eta))/D-\nu^3\xi^6
\]
となる。実際 $cd=w+\nu a^2$ で $d$ を消去でき、
$B'$ の $c$ 局所化は $k[a,b,w]$ 上のskew Laurent表示であり、$v^3=F$ は
元の三次切断に等しい。$d$ のchartも、$c,d$ と $C,D$ を交換し
$q$ を $q^{-1}$ に替えれば同じ形になる。$\nu$ は変わらない。

$F=0$ は原点を通らない。$\xi=0$ 上では $m=2$ なら零点がなく、
$m=3$ なら零点は $\eta\ne0$ で単純である。
$\xi\ne0$ 上では $z=\eta/\xi$、$x=\xi^{-3}$ とおくと、
$F=0$ は $g(x,z)=0$ のエタール引き戻しである。
従ってこの分岐曲線が滑らかであることは $\Delta$ が重根を持たないことに等しい。
$g$ の定数項が非零なので $g=g_x=g_z=0$ では $x\ne0$ となり、
$g_x=0$ に代入することで重根との同値が分かる。

局所大域次元は可換不変環 $D_0=k[\xi,\eta]^{\langle\tau\rangle}$ 上で調べる。
原点以外の閉点の完備化では、係数環が三つの等しい二次元正則局所環の積へ分かれ、
$v$ は三冪等元を巡回させる。
$v^3=F$ が単元なら行列環になり、$F=0$ の点ではこの環は
\[
 \bigl(S[t]/(t^3-F)\bigr)\rtimes\mu_3
\]
と同型である。ここで $S$ は不変環の完備局所環で、
$\mu_3$ は $t$ を三乗根倍する。
同型は三つの冪等元を群環の指標冪等元に、$v$ を $t$ に送ることで得られる。
$F$ の微分が非零なら $S[t]/(t^3-F)$ は正則局所環であり、
有限群環は大域次元2を持つ。これは誘導・制限と標数零での平均化で従う。
微分が零なら可換超曲面の剰余体の射影次元は無限である。
その誘導加群は三つの指標単純の和なので、
少なくとも一つは群環上で無限射影次元を持つ。
この単純加群は完備化前の中心ファイバーから来る。
従ってその点があれば元のchartも有限大域次元を持たない。
H3のtailsから中央開部分への局所化は高次Ext消滅を保つので
（補題\ref{lem:r13-one-chart-smoothness}の局所化の議論）、
H3ではない。

原点の完備化では $F$ は単元であり、中央不変な三乗根で $v$ を割ると
$k[[\xi,\eta]]\rtimes\mu_3$ を得る。従ってここも有限大域次元である。
全閉点の完備化での有限射影次元の一様上界は、中心有限性と完備化の忠実平坦性により、
有限生成加群の有限自由解消のsyzygyの射影性に降下する。

\paragraph{残る局所chartとH3認識}
中央三次 $a^3,b^3,X,Y$ の開集合は射影中心を覆う。
既に扱った $c,d$ の両chartの外側では $X=Y=0$ であり、
まず $a\ne0$ の部分を扱う。
$a=1$ の次数零環を中心 $k[z,x]$ の $X=Y=0$ に対応する点で完備化すると、
$w^3=-\nu^3$ が剰余関係なので $w$ の三固有冪等元が分裂する。
$c,d$ は三冪等元間の逆向きの矢を与え、
その二つの積は $cd=w+\nu$ と $dc=qw+\nu$ である。
三つの隣接辺のうち二つではこの積が単元であり、
二辺だけで全冪等元が互いに同値になる。
完備化後も中心上階数9の自由加群なので、この環は中心完備局所環上の三次行列環である。
ここで $X,Y,w$ を $a=1$ によって次数零の元と同定した。

$a=X=Y=0$ となる射影点は $m=2$ の場合の $b$ 軸だけである。
$b=1$ として $a,c,d$ を次数零元で書く。
$P(a,1)+CX+DY=0$ の $a$ の一次係数は非零なので、
中心をこの点で完備化すると $a=f(X,Y)\in(X,Y)k[[X,Y]]$ と一意に解ける。
完備環 $A$ は $k[[X,Y]]$ 上 $c^i d^j$ ($0\le i,j<3$) を基底とし、
\[
 c^3=X,\qquad d^3=Y,\qquad
 dc=qcd+(1-q)\nu f(X,Y)^2 .
\]
重み $(3,3,1,1)$ を $X,Y,c,d$ に与えた降下濾過で、
補正項は重み6以上である。
自由基底より、随伴次数環は量子平面である。
この濾過は完備である。正重みのイデアルを $J=(c,d)$ とすると、
$A/J=k$ であり、正重み元の幾何級数から $J$ は $A$ のJacobson radicalである。
任意の重み $n>0$ の元は $c,d$ と重み $n-1$ 以上の元との積の和に書ける。
従って右加群写像 $(u,v)\mapsto cu+dv$ はstrictである。
先頭syzygy $(-qd,c)$ は、その像が重み6以上なので
重み5以上の補正を加えて実際のsyzygyに持ち上がる。
量子平面の解消と完備性により、この一つのsyzygyが核を自由に生成する。
具体的には先頭項が定める倍数を引き、重みが上がるごとにこれを反復する。
左側も同様である。従って $k$ は左右とも射影次元2を持つ。
有限生成加群の極小自由解消を $k$ とテンソルすると微分が消えるため、
一般の有限生成加群の射影次元も高々2である。
これで $b$ 軸の完備局所環も有限大域次元を持つ。

従って $\Delta$ が平方自由なら、三次Veroneseの中心開被覆による
局所Extとそのコホモロジーから $\qgr R'$ は有限大域次元を持つ。
具体的には中心有限の代数の層として、局所Extの次数は高々2、
射影中心のコホモロジー次数も高々2である。
Zhang同値により $\qgr R$ も有限大域次元である。
$R$ はCY AS4の正則中心三次商だからAS Gorenstein次元3、parameter 1である。
有限大域次元と双対性から $S=(-1)[2]$、従って大域次元は2となる。
中心有限整域のgeneric商は斜体の有限次元加群圏なので、構造対象は
Gabriel--Krull rank-one torsion-freeである。
ここでは $R$ と $R'$ の三次Veroneseは同じであり、中央有限性はこの共通環でも成り立つ。
深さ3から飽和と例外性、$\chi$ と一次生成からAZ豊富性が従う。
Hilbert値および二次包絡は正則三次商であることから従い、H3の条件が揃う。

最後に $P=b^3,C=D=1$ では $\Delta=z^6-4\nu^3$ は平方自由である。
元の包絡は一次正規元が張る空間が $\langle a,b\rangle$ に限られるので非skewである。
実際、$(a,b)$ を法とした量子平面で一次正規元は $c,d$ の座標軸に限られる。
$c+u$、$u\in\langle a,b\rangle$ が正規なら $d(c+u)\in(c+u)B_1$ である。
順序語の $cd,c^2,ac,bc$ を比較すると右辺は $(c+u)q^2d$ しかあり得ず、
元関係の $\kappa a^2$ が残って矛盾する。
$d+u$ も $c$ との関係を比較して同様である。
従って四つの独立な一次正規生成元は存在しない。
\end{proof}

\begin{corollary}\label{cor:cubic-CY-one-segre-square-split-resonance-groupoid}
原始三乗根の一方を固定して $Z(s)$ のH3切断を分類する。
$C=D=1$ に規格化した表示の全幾何学的群oidは
\[
 \left[
 \{P\in k[a,b]_3:\deg_zP(1,z)\ge2, P(1,z)^2-4\nu^3\text{ が平方自由}\}
 \big/
 (\operatorname{Aff}_1\times\mu_6)
 \right]
\]
である。
作用の条件は、二つの表示間に
\[
 P'(a,b)=\mu^{-3}P(a,u a+v b),\qquad
 u\in k,\ v\in k^*,\ \mu^6=1
\]
を満たすパラメータが存在することであり、各パラメータが一つの矢を与える。
一般の表示の慣性は $\mu_3$、この群oidの次元は2である。
\end{corollary}
\begin{proof}
$U_0$ の二次元ブロックの組は内在的に復元できる。
実際、$U_0=L_1L_2$ 内の分解可能な非零二次式の一次因子の集合は
二つのブロックの和集合である。これは二次式の両純ブロックへの射影が零であることから従う。
$U/U_0$ の各ブロックの階数が1と2なので、同型は順序も保つ。
従って $a$ の直線と $kc,kd$ の組を保つ。
固定した $q$ において $c,d$ を交換すると $L_1$ との交換係数が逆になり、
同型にならない。別の原始三乗根値の表示へは、この交換が同型を与える。
実際 $s\mapsto-s$ で $q\mapsto q^{-1}$、
係数 $\kappa=8s/(2s-1)^2$ は $-q^{-2}\kappa$ に移る。

従って固定包絡の次数付き自己同型は
\[
 (a,b,c,d)\longmapsto (r a,u a+v b,\mu c,\lambda d),
 \qquad r,v,\mu,\lambda\in k^*,\quad \mu\lambda=r^2
\]
に尽きる。最後の条件は $dc-q^2cd=\kappa a^2$ によるもので、
逆にこれらの条件で全関係が保たれる。
H3の同値は反Serre環の次数付き同型なので、二次包絡へ一意に持ち上がる。
点付き自然同型で割ると共通の生成元スカラーが消え、$r=1$ と規格化できる。

$CD\ne0$ の切断直線は $\mu\lambda=1$ の作用と三次式の共通尺度で
$C=D=1$ に移せる。二つの規格化済み表示の間では $\mu^6=1$ が必要十分で、
共通尺度が $\mu^{-3}$ になる。従って上の作用を得る。
一般の三次式を平行移動により $L z^3+M z+N$ と書き、
$LMN\ne0$ とすると、これを符号倍まで保つアフィン変換は恒等だけである。
実際、二次係数で平行移動が消え、一次と定数の比較で倍率も符号も1となる。
よって一般慣性は $\mu^3=1$ であり、群oidの次元は $4-2=2$ である。
\end{proof}

\begin{proposition}\label{prop:cubic-CY-one-segre-square-split-boundary-cross}
等方外側の残る $b_5(s,\theta)$、
下行片成分の両交差型 $L^\pm_{w,e}$、および
両下行・$cd$ 係数零の交差型 $Z_{s,h}$ にはH3切断が存在しない。
従ってこの平方・積追加射影でなお未判定の包絡は、
固定出典の $M_{s,w,z}$ 一族に限られる。
\end{proposition}
\begin{proof}
$b_5$ の関係を順序 $a<b<c<d$ で書き換える。
中心行列から、四つの可換二元三次列が零の行だけを取ると、
対応する16列の小行列式として
\[
 -\frac{17179869184,s^4(48s^2+1)^2}{(4s-1)^{14}(4s+1)^9}
\]
を得る。逆順語の書換え則の分母は許容単元だけである。
残る $48s^2+1=0$ では同じ部分行列を再計算し、
$\theta$ に依存しない非零16次小行列式を得る。
よって全域で中心三次の非二元座標は零であり、二元三次の分解によって排除する。

次に固定出典\cite{AS4CYpinned}の下行片成分と両下行・$cd$ 係数零の必要表示を用いる。
片成分の二符号は同じ係数の反対代数なので中心三次の消滅は一方を示せばよい。
元の六二次関係の次数2--4の射影と中心20次小行列式から、
片成分の許容域内では次の五曲線だけが残る。
\[
\begin{gathered}
 w=-32,\qquad e w(w+16)=1,\qquad e w(w+32)=1,\\
 e^2w^3(w+16)-2e w(w+16)+1=0,\\
 e^2w^3(w+16)^2-2ew(w^2+24w+256)+w+16=0 .
\end{gathered}
\]
第四曲線は $t=ew(w+16)$ を取ると
$w(t-1)^2=16(2t-1)$ となり、$w(w+16)(t-1)\ne0$ なので式から $w,e$ を復元できる。
第五曲線では $w=-16$ は $e\ne0$ と両立しない。
\[
 y=\frac{e w^2(w+16)^2-w^2-24w-256}{8}
\]
とおくと $y^2=5w^2+128w+1024$ であり、
$t=(y-32)/w$ から $w=(128-64t)/(t^2-5)$ を復元する。
$t^2=5$ は元の式と矛盾するので、この有理化に漏れはない。

両下行・$cd$ 係数零の型では $b_0=s^2-1/16$、
$x=b_0h/s$ と置く。元の許容単元を除く同じ因子計算から得る五つの残余は
\[
\begin{gathered}
 s^2(x+2)^2(x+4)=(3x+4)/4,\qquad
 s^2(x+2)^2=(x+1)/4,\\
 x=-2,\qquad (s+1/4)h+2=0,\qquad
 s^2(x+2)^2-sx/3-1/4=0 .
\end{gathered}
\]
第一、第二、第五は $u=2s(x+2)$ として
\[
 x=\frac{4(1-u^2)}{u^2-3},\qquad
 x=u^2-1,\qquad
 x=-\frac{2(u^2-1)}{u^2-1-2u/3}
\]
で有理化される。最初の二式では $x=-2$ は元の条件に反する。
第五式の $x=-2$ は $(s,h)=(3/8,-48/5)$ を与え、別に計算する。
それ以外は分母を元の式に代入することで非零性を得る。
第三と第四の曲線は直接 $h$ について解ける。

この各有理曲線上で商射影と中心行列を元関係から再計算する。
更に残る全一次・二次既約因子の数体と、第五式の別点から商を取り直すと、
中心三次行列は常に階数20であった。
\texttt{verify\_cubic\_CY\_one\_segre\_square\_split\_lower.py} と
対応する証明書はこの全因子と再計算を保持する。
\end{proof}

\begin{proposition}\label{prop:cubic-CY-one-segre-square-split-open-Wminus}
残る $M_{s,w,z}$ のうち
$(4s-1)^2w+64=0$ または $(4s+1)^2w+64=0$ の切断はH3でない。
\end{proposition}
\begin{proof}
固定出典\cite{AS4CYpinned}の $M_{s,w,z}$ の必要表示と、
反対代数が $s\mapsto-s$ に対応する式 (LO11) を使う。
反対代数で中心は変わらないので $w=-64/(4s-1)^2$ だけを調べる。
元の六関係にこの値を代入し、生成元順を $(b,a,c,d)$ とした逐次商で
次数2から4の射影と中心20次小行列式を取る。
分母と小行列式の分子・分母の既約因子のうち、
元のパラメータ域で単元でないものは次の三式だけである。
\begin{align*}
 &sz+\tfrac3{32}s^5-\tfrac5{128}s^4-\tfrac1{256}s^3
     +\tfrac3{1024}s^2-\tfrac1{8192}s-\tfrac1{32768},\\
 &sz+\tfrac1{16}s^5-\tfrac1{64}s^4-\tfrac1{128}s^3
     +\tfrac1{512}s^2+\tfrac1{4096}s-\tfrac1{16384},\\
 &sz+\tfrac18s^5-\tfrac5{64}s^4+\tfrac1{64}s^3
     -\tfrac1{512}s^2+\tfrac1{2048}s-\tfrac1{16384}.
\end{align*}
各式で $s\ne0$ なので $z$ を解き、元六関係から次数2から4の全射影を取り直す。
各曲線では再計算した小行列式と全射影の分母の因子は許容単元だけであり、
中心行列の階数は20である。
この計算は \texttt{verify\_cubic\_CY\_one\_segre\_square\_split\_lower.py}
の最後のチャートで再現できる。
従ってこの二境界では全中心三次が零である。
\end{proof}

\begin{proposition}\label{prop:cubic-CY-one-segre-square-split-open-Mzero}
$M_{s,w,z}$ の $(16s^2+1)w+64=0$ の範囲では、
中心三次が零である。従ってこの族に残るH3候補には
$(16s^2+1)w+64$ と $(4s\pm1)^2w+64$ の三因子の非消失が必要である。
\end{proposition}
\begin{proof}
この境界では $16s^2+1\ne0$ であり、$w=-64/(16s^2+1)$ と置ける。
生成元順 $(a,c,b,d)$ で逐次商を取り、次数2から4の全射影と中心20次小行列式の
非単元因子を計算すると、次の三直線型の式と一二次式が得られる。
\begin{align*}
 &z-\tfrac1{64}s^2-\tfrac1{1024},\\
 &sz+\tfrac1{16}s^5-\tfrac1{64}s^4-\tfrac1{128}s^3
      -\tfrac1{512}s^2-\tfrac3{4096}s-\tfrac1{16384},\\
 &s^2z+\tfrac1{16}s^6-\tfrac1{256}s^4-\tfrac1{4096}s^2+\tfrac1{65536}.
\end{align*}
二次式は $s^2z^2+B(s)z+C(s)$ であり、
\[
 B(s)=s^6/8-s^4/64-3s^2/2048,\qquad
 B(s)^2-4s^2C(s)=-s^2(16s^2+1)^4/2^{26}.
\]
従って二次式の零点は $\mathbb Q(i)(s)$ 上で二つの有理グラフに分かれ、
その一方は
\[
 z=\frac{-B(s)+i\,s(16s^2+1)^2/8192}{2s^2}
\]
である。他方は $\mathbb Q(i)$ の共役である。

三つの一次式ではそれぞれ $z$ を解いて元関係から再計算すると、
中心行列の階数20が全許容値へ特殊化する。
二次式の指定グラフでも射影と小行列式の非単元を
$\mathbb Q(i)[s]$ で因数分解すると、許容な残余は高々 $s=1/2-i/4$ の一点となる。
この点も $\mathbb Q(i)$ で階数20である。
係数の共役を取れば他方のグラフも全域で階数20となる。
以上は同じ公開検証器と証明書で再現できる。
\end{proof}

\begin{proposition}\label{prop:cubic-CY-one-segre-square-split-M-quadratic}
$M_{s,w,z}$ の許容域(LO4)全体で
\[
 U=a-(c+kd)/s,\qquad V=s(b-d)
\]
は可換であり、次の二次式は正則かつ中央である。
\[
 Q_0=\left(\frac{sw}{4(kw+1)}\right)^2 U^2-2sm\,UV+V^2,qquad
 m=\frac{(32k+1)w+32}{16s(kw+1)} .
\]
従って $Z(M_{s,w,z})_3$ は $Z(M_{s,w,z}/(Q_0))_3$ に単射し、
二次商の次数二、三、四の次元は $9,16,25$ である。
\end{proposition}
\begin{proof}
$[U,V]$ は(LO1)--(LO2)の実交換子である。
(LO5)の関係を、混合四列を二次係数の一次式のまま保つ基底に変更する。
$(N,b,D)$ を(LS4)の元係数とし、$D=C+\gamma J-N^{\mathsf T}A^{-1}N$ とする。
四列 $(U_{\rm sym}+KN^{\mathsf T}A^{-1})$ に右から $A$ を掛け、
追加列の $D$ を含む項をこれらの列の線形結合で置き換える。
従って新たな六列は
\[
 U_{\rm sym}A+KN^{\mathsf T},\quad Kb,\quad
 a^2+c\circ d+
 \frac{U_{\rm sym}N_{*,3}+K(C+\gamma J)_{*,3}}{b_3} .
\]
ここで $b_3=1$、$c\circ d=(cd+dc)/2$ である。
これらの六列から次数三商を構成すると、三変数有理関数体上で
$Q_0$ の全四交換子が厳密に零となる。
係数の分母は(LO4)の単元で、次数三商の次元は(LO4)全域で一定なので、
中央性は全許容値へ特殊化する。
計算は \path{verify_cubic_CY_one_segre_square_split_M.py} で再現する。

$U,V$ は独立で、二次式の $V^2$ 係数は1である。
可換二元二次としての一次因子は整域でいずれも非零なので $Q_0$ は正則である。
実交換子空間の次元は一だから一次中心は零であり、
命題\ref{prop:cubic-CY-one-segre-double-split-quadratic-reduction} と同じ
正則列の議論で単射とHilbert値が従う。
\end{proof}

残る三因子非零域では、次の座標が中央二次商を扱いやすくする。
\[
 X=4s,\quad T=\frac{kw-1}{kw+1},\quad Z=\frac{wz-k}{k},\qquad
 w=\frac{1+T}{k(1-T)},\quad z=\frac{k(Z+1)}{w}.
\]
(LO4)の非消失条件は、この座標で
\[
 X(X^2-1)(T^2-1)
 \prod_{\epsilon=\pm1}
 \bigl(Z(X+\epsilon T)^2-T(X^2+1)-2\epsilon X\bigr)\ne0 .
\]
既に排除した三境界は $X+T=0$, $X-T=0$, $X^2+T=0$ である。
(LO6)は $(X,T,Z)\mapsto(-X,-T,-Z)$、
反対代数(LO11)は $(X,T,Z)\mapsto(-X,T,Z)$ となる。
中央二次の $UV$ 係数は $-4(X^2+T)/(X^2-1)$、
二元二次式としての判別式は $16(X^2-T^2)/(X^2-1)$ である。

\begin{corollary}\label{cor:cubic-CY-one-segre-square-split-M-generic-zero}
$M_{s,w,z}$ とその中央二次商では、稠密開集合で中央三次は零である。
\end{corollary}
\begin{proof}
$(s,w,z)=(1,1,1)$ は(LO4)を満たす。
この点で元六関係の中心行列は階数20、中央二次商は階数16である。
\path{verify_cubic_CY_one_segre_last_generics.py} で両方を厳密に再計算する。
次数三、四の商の次元は全許容値で一定なので、
階数の下半連続性からそれぞれ一般点でも同じ階数となる。
\end{proof}
